% Generated draft: requires mathematical and source review.
\documentclass[11pt]{article}
\usepackage[margin=25mm]{geometry}
\usepackage{fontspec}
\setmainfont{Latin Modern Roman}
\usepackage{amsmath,amssymb,mathtools,unicode-math}
\setmathfont{Latin Modern Math}
\usepackage{enumitem,needspace,xurl,hyperref}
\setlength{\parindent}{0pt}
\setlength{\parskip}{5pt}
\setlength{\emergencystretch}{3em}
\newcommand{\N}{\mathbb{N}}
\newcommand{\Z}{\mathbb{Z}}
\newcommand{\Q}{\mathbb{Q}}
\newcommand{\R}{\mathbb{R}}
\newcommand{\C}{\mathbb{C}}
\newcommand{\F}{\mathbb{F}}
\newcommand{\A}{\mathbb{A}}
\newcommand{\PP}{\mathbb{P}}
\newcommand{\E}{\mathbb{E}}
\newcommand{\eps}{\varepsilon}
\newcommand{\dd}{\,\mathrm{d}}
\begin{document}
\section*{UnsolvedMath ID: 30001012}
\textbf{Source problem code:} catalogue entry
\subsection{Definitions and mathematical statement}
Let \(H\cong\ell^2(\N)\) be a separable infinite-dimensional complex Hilbert space, and let \(B(H)\) be its algebra of bounded operators. Let \(C\subseteq B(H)\) be a unital maximal abelian self-adjoint subalgebra, so
\[
 C'\cap B(H)=C,
\]
isomorphic to \(L^\infty((0,1),\dd t)\). Thus the MASA is diffuse, with no nonzero minimal projections; it is not the atomic diagonal MASA.

Let \(F_2\) be the free group on two generators and \(A=C_r^*(F_2)\) its reduced group \(C^*\)-algebra, the operator-norm closure of the left regular group algebra on \(\ell^2(F_2)\). Form
\[
 B=B(H)\otimes_{\max}A,
\]
the maximal \(C^*\)-completion of the algebraic tensor product. Define the relative commutant by
\[
 (C\otimes1)^c=\{x\in B:x(c\otimes1)=(c\otimes1)x\text{ for all }c\in C\}.
\]
The right-hand tensor product in the question means the norm-closed canonical image
\[
 C\otimes A=\overline{C\odot A}^{\,\|\cdot\|_B}\ \subseteq B.
\]
Since \(C\) is abelian, hence nuclear, its minimal and maximal tensor norms with \(A\) agree. The canonical map into \(B\) is injective: composing with \(B(H)\otimes_{\max}A\to B(H)\otimes_{\min}A\) gives the injective minimal-tensor inclusion. This makes the notation on the right unambiguous.

The question is whether
\[
 (C\otimes1)^c\stackrel{?}{=}C\otimes C_r^*(F_2).
\]
The right side is always contained in the left because \(C\) is abelian. The substantive assertion is that every element commuting with all \(c\otimes1\) in the maximal ambient tensor product is a norm limit of finite sums \(\sum_i c_i\otimes a_i\), with \(c_i\in C\) and \(a_i\in C_r^*(F_2)\). It does not ask for a weak-operator closure, a von Neumann tensor product, or the analogous equality in a minimal ambient tensor product.

The supplied assessment dated 2026-08-21 attributes this exact diffuse-MASA relative-commutant question to Wassermann and reports it as open, with no general resolution located by the supplied curation.
\subsection{Short English statement}
Inside the maximal tensor product of \(B(H)\) with the reduced free-group algebra, does commuting with the diffuse MASA force an element to lie in the norm-closed tensor product with that MASA? The reverse inclusion is the entire open assertion.
\subsection{Sources}
\begin{itemize}
\item Local UnsolvedMath registry, permanent ID \(30001012\), exact source code \texttt{catalogue entry}, proposed year 2008; no individual proposer is supplied, while the assessment attributes the question to Wassermann.
\item \emph{\(C^*\)-Algebras (2009)}, Oberwolfach Reports, DOI 10.4171/owr/2008/37, pp. 2095--2164. \url{https://doi.org/10.4171/owr/2008/37}. The supplied title contains escaping artifacts; the mathematical title is rendered here in standard notation.
\item Wassermann, \emph{Tensor products of maximal abelian subalgebras}, \url{https://arxiv.org/abs/0711.0131}.
\item The supplied \texttt{corrected\_verified} note identifies a complete relative-commutant question with restored maximal-tensor notation. The open label and assessment checked 2026-08-21 are supplied dated curation, not independently verified; exact current status requires review.
\end{itemize}
\end{document}
